\section{Successors, fundamental sequences and cofinality}
\label{sec:fundseq}

Write $P\mathbin{\|}Q$ for list concatenation, $[b]$ for a one-entry
list, and $[b]^k$ for $k$ copies of $b$ in a list.
All arguments of ordinal comparisons in this section are canonical terms.
We use the convention that a limit ordinal is positive.
The root of even the empty term counts as one node. Thus its node count is
\[
 v([a_1,\ldots,a_k])=1+\sum_{i=1}^k v(a_i).
\]

\begin{proposition}[Successor and finite remainder]
The immediate successor of $x$ is $S(x)=x\mathbin{\|}[0]$.
A nonempty term is a successor exactly when its last exponent is zero.
Every term has a unique expression $S^r(L)$ with $r$ a natural number and
$L$ zero or a positive term ending in a positive exponent.
\end{proposition}
\begin{proof}
A proper extension of $x$ begins with an additional exponent at least zero,
so it is at least $S(x)$. If a term exceeds $x$ by a first unequal
exponent, that same difference puts it above $S(x)$.
Thus $x<y$ is equivalent to $S(x)\le y$.
Deleting a last zero inverts $S$. Counting and deleting all trailing
zeros gives $r,L$ and reversing this operation recovers the input.
Maximality of the trailing zero run proves uniqueness. In particular the
finite terms are the lists $[0]^r$. The terms below $\omega=[1]$, where
$1=[0]$, are exactly these lists, because every positive exponent is at least
$1$. Thus $\omega$ is the first infinite term.
\end{proof}

For a term $\lambda=P\mathbin{\|}[b]$ ending in $b>0$, define
\[
 F(\lambda,n)=
 \begin{cases}
 P\mathbin{\|}[c]^{n+1},& b=S(c),\\
 P\mathbin{\|}[F(b,n)],& b\text{ ends in a positive exponent}.
 \end{cases}
\]
Give $F$ the default value zero on other inputs. The computational call
in the second line follows the last child in the input tree.
The number of calls is bounded by $v(\lambda)$. The two clauses give
$v(F(\lambda,n))\le(n+1)v(\lambda)$ by addition of node counts along
this path. One can therefore implement the definition by bounded
descent with a stack of saved prefixes and a reverse traversal for
reassembly. For instance,
\[
 F(\omega,n)=n+1,\qquad
 F(\omega^2,n)=\omega(n+1),\qquad
 F(\omega^\omega,n)=\omega^{n+1}.
\]
Here numerals in ordinal expressions stand for the corresponding finite
terms. The indexing convention starts each example strictly above zero.

\begin{theorem}[A quantitative cofinality theorem]
\label{thm:fsquant}
If $\lambda$ ends in a positive exponent, $F(\lambda,n)$ is canonical,
and for every $n$,
\[
 F(\lambda,n)<F(\lambda,n+1)<\lambda.
\]
More precisely, whenever $y<\lambda$ and $n\ge v(y)$, one has
$y<F(\lambda,n)$.
\end{theorem}
\begin{proof}
In the successor-exponent clause, every entry of $P$ is at least $b>c$.
The extra copy of $c$ gives strict increase by proper extension. The
first replacement of $b$ by $c$ places both results below $\lambda$.
In the limit-exponent clause the same conclusions hold by applying the
comparisons for the smaller tree $b$ to the only varying entry. This
also justifies canonicality at every return from the recursive call.

For the quantitative claim compare $y$ with $P\mathbin{\|}[b]$. If
$y$ ends in or at $P$ as a prefix, the result has a nonempty suffix
after $P$ and exceeds $y$. If $y$ first differs inside $P$, the
smaller entry of $y$ still decides the comparison. The remaining case is
$y=P\mathbin{\|}[d]\mathbin{\|}Q$ with $d<b$. If $b=S(c)$ then
$d\le c$. For $d<c$ the first new entry decides. For $d=c$ all
entries of $[d]\mathbin{\|}Q$ are at most $c$, and this list has
at most $v(y)$ entries. A list of $n+1$ copies of $c$ therefore
exceeds it at a first smaller entry or by proper extension. If $b$ ends
in a positive exponent, apply the claim to the smaller input pair
$(b,d)$, since $v(d)<v(y)\le n$. It gives $d<F(b,n)$ as required.

These verifications use ordinary induction on the length of the
calculation. Explicitly, for fixed inputs store the visited pairs
$(\lambda,y),(b,d),\ldots$ and the common $n$. At each pair the local
comparison either decides the result outright or invokes the next pair.
Check the implications backwards with a bounded conjunction over the
visited nodes. For strict increase use the analogous list of rightmost
ancestors with parameters $n,n+1$. This makes the induction predicate
a primitive recursive Boolean test.
\end{proof}

\begin{corollary}[The exact first cofinality index]
\label{thm:exact-fs}
Let $M(y)$ be the largest length of a run of identical consecutive
children anywhere in the tree $y$, with empty maximum zero.
For each positive limit $\lambda$ and $y<\lambda$ a primitive recursive
function $m(\lambda,y)$ satisfies
\[
   y<F(\lambda,n)\quad\Longleftrightarrow\quad n\ge m(\lambda,y),
   \qquad m(\lambda,y)\le M(y).
\]
In particular $M(y)$ can replace $v(y)$ in
Theorem~\ref{thm:fsquant}. For every natural $k$, equality holds in
the bound when $\lambda=\omega$ and $y=k$.
\end{corollary}
\begin{proof}
Write $\lambda=P\mathbin{\|}[b]$. If $y$ is a prefix of $P$,
including equality, or first differs from $\lambda$ inside $P$,
return zero. The whole prefix $P$ survives and is followed by at
least one entry, so the desired strict comparison holds at every
index. Otherwise $y=P\mathbin{\|}[d]\mathbin{\|}Q$ with $d<b$.
If $b=S(c)$ and $d<c$, return zero again. If $d=c$, return the length
$r$ of the initial run of $c$ in $[d]\mathbin{\|}Q$. The new suffix
of $F(\lambda,n)$ has $n+1$ copies of $c$. For $n+1\le r$ this suffix
is a prefix of the suffix of $y$, possibly equal. For $n+1>r$ it
exceeds that suffix because $y$ ends after the run or continues with
a smaller entry. This proves the threshold $n\ge r$, and $r\le M(y)$.
If $b$ is a limit, return $m(b,d)$. The new suffix has the single
entry $F(b,n)$, so $y<F(\lambda,n)$ is equivalent to $d<F(b,n)$.
In particular equality at that entry cannot give strict inequality,
even when $Q$ is empty. The recursion descends into proper subtrees, and
$M(d)\le M(y)$ as $d$ is a child of $y$. A bounded stack of these
input pairs and a backwards scan prove both assertions by the same
quantifier-free induction as before. Computing $M$ itself requires
only counting runs at each node and taking their finite maximum.
Finally $F(\omega,n)=n+1$ and $M(k)=k$, so $m(\omega,k)=k$.
\end{proof}

\begin{example}
Let $\lambda=\omega^{\omega^2}$ and $y=\omega^{\omega\cdot3+7}+42$.
Comparison of the leading exponents reduces the first-index problem to
$\omega^2$ and $\omega\cdot3+7$. The latter begins with three copies of
the exponent $1$, so $m(\lambda,y)=3$. Indeed
$F(\lambda,2)=\omega^{\omega\cdot3}<y<F(\lambda,3)=\omega^{\omega\cdot4}$.
The final $42$ units do not enter the calculation.
\end{example}

\begin{corollary}[Characterizations and cofinalities]
For $x>0$ the following four conditions are equivalent.
\begin{enumerate}[nosep,label=(\roman*)]
\item The last exponent of $x$ is positive.
\item The term $x$ has no predecessor.
\item $S(y)<x$ for every $y<x$.
\item The strict predecessors of $x$ have least upper bound $x$.
\end{enumerate}
The cofinalities of zero, a successor and a positive limit in this
presentation are respectively $0,1,\omega$. Here a cofinal list or
sequence below $x$ must supply an entry $\ge y$ for each $y<x$.
\end{corollary}
\begin{proof}
For a positive last exponent, take $n=v(y)$ in
Theorem~\ref{thm:fsquant}. Then $S(y)\le F(x,n)<x$.
Every proposed upper bound $y<x$ of the predecessors is exceeded by
$F(x,n)$. A final zero, on the other hand, exhibits a predecessor $p$
which is itself the greatest strict predecessor and their least upper
bound. This separates the two cases constructively.

The empty list suffices at zero and the singleton containing $p$ suffices
at $S(p)$, where an empty list fails on $p$. At a positive limit the
displayed $F$ gives an increasing sequence. It exceeds the maximum of
each nonempty finite list of predecessors at the explicit index given
above. The empty list misses zero. Since the terms below $\omega$ are
the finite ones, this also proves minimality of the stated cofinality.
In particular $\omega$ is the least positive limit. The three values give
$\operatorname{cf}(\operatorname{cf}(x))=\operatorname{cf}(x)$ by
inspection.
\end{proof}

\begin{proposition}[A minimum problem already at $\omega+1$]
Let $b(n)$ be any specified primitive recursive Boolean predicate,
possibly with numerical parameters. In $[0,\omega]$ form the decidable
inhabited predicate $A_b$ consisting of $\omega$ and exactly those
finite $n$ for which $b(n)=1$. Then $A_b$ has a least member if and
only if
\[
       (\exists n\ b(n)=1)\ \lor\ (\forall n\ b(n)=0).
\]
\end{proposition}
\begin{proof}
Given a least member, decide whether it is finite or equals $\omega$.
A finite least member supplies a positive instance. If the least
member is $\omega$, any $n$ with $b(n)=1$ would be a smaller member.
Boolean decidability now gives $b(n)=0$ at each input. Conversely,
from a witness $n$ scan the $n+1$ finite positions up to it to find
the first positive one. This is least among all finite members and
lies below $\omega$. In the all-zero case $\omega$ itself is least.
Thus the assertion that all such inhabited decidable subsets have a
minimum already entails the limited principle of omniscience for
primitive recursive predicates, understood as a first-order schema.
\end{proof}

\subsection{The tower endpoint and continuity of the principal-power map}
Put $p(x)=[x]$, the notation for $\omega^x$. Singleton comparison gives
$p(x)<p(y)$ exactly when $x<y$. In fact $x<p(x)$ for every finite
term. This holds at zero. If $a$ is the first child of $x$, then
$a<p(a)$ by the same assertion for the smaller tree, and $p(a)\le x$
by the prefix comparison. Thus $a<x$, whence $x<[x]$ by comparison
of the first entries. This argument is a bounded backwards check on the
leftmost branch. At a positive limit $\lambda$, the sequence
$p(F(\lambda,n))=F(p(\lambda),n)$ is cofinal below $p(\lambda)$.
In particular, given $z<p(\lambda)$, the bound $n=v(z)$ supplies an
argument $u=F(\lambda,n)<\lambda$ with $z<p(u)$. This proves
continuity in the increasing, least-upper-bound sense.

Define $T_0=1$, $T_{n+1}=p(T_n)$, and let $h(0)=0$,
$h(x)=1+\max_i h(a_i)$ for a nonempty term with children $a_i$.
The tower is strictly increasing. In fact,
\[
                 x<T_n\quad\Longleftrightarrow\quad h(x)\le n.
\]
To see this, first prove $x<T_{h(x)}$, and, for $h(x)=d>0$,
$T_{d-1}\le x$. At height zero the upper inequality is $0<1$.
At positive height $d$, all children are below $T_{d-1}$ by their
upper inequalities and monotonicity of the tower. Comparing first
entries gives $x<[T_{d-1}]=T_d$. When $d=1$, nonemptiness gives
$x\ge1=T_0$. When $d>1$, a child of height $d-1$ is at least
$T_{d-2}$ by its lower inequality. The first child is at least this
large and therefore $x\ge[T_{d-2}]=T_{d-1}$. Both estimates can be
verified on the finite table of nodes of the given tree, in postorder.
Together they give the asserted equivalence.

Thus the towers are cofinal in the entire term order, with witness
$h(x)$ for a strict upper bound of $x$. There is also an exact
fixed-point formulation using a one-symbol extension. Adjoin a last
point $E$ and set $p(E)=E$. Then $E$ is the least fixed point, since
$x<p(x)$ for every ordinary term. The towers have supremum $E$, and
$p$ is continuous there since, for $x<E$, $p(T_{h(x)})>x$. This models the
first epsilon fixed point at the endpoint of this presentation. The
unextended term order itself has no largest element.

\section{Order topology by explicit neighborhood tests}
\label{sec:ordtop}
All statements about a family of points below refer to its specified
membership predicate. A basic neighborhood of a nonzero point $x$ can
be shrunk to $(b,S(x))$ with $b<x$, since the right-hand endpoint can
be moved to $S(x)$ and zero supplies a possible lower endpoint.
At zero use $[0,1)$. In an initial segment intersect these intervals
with the segment. A point $x>0$ is therefore a limit point of a
predicate $A$ precisely when for every $b<x$ there is a $y$ with
$b<y<x$ and $A(y)$.

Successors $S(p)$ are isolated by $(p,S(S(p)))$, and zero is isolated.
At a positive limit $x$ the sequence $F(x,n)$ gives points approaching
$x$ strictly from below. The intervals $(F(x,n),S(x))$ form a
decreasing countable local basis by Theorem~\ref{thm:fsquant}. As a
result, an initial segment $[0,a)$ is discrete exactly when
$a\le\omega$. For $a\le\omega$ every point is finite. For $a>\omega$
the segment contains $\omega$ and all the finite points approaching it.
On $[0,\omega]$, a fixed decidable membership predicate $U$ describes
an open set exactly when it omits $\omega$ or it contains $\omega$ and
there is an $N$ such that $U(n)$ holds for every finite $n\ge N$.
Indeed, the isolated points impose no condition. A neighborhood of
$\omega$ contains a final tail and every such tail together with
$\omega$ contains a basic neighborhood. Thus in the case containing
$\omega$ the complement can be given by a finite list, obtained by
checking the $N$ positions below the cutoff.

\begin{theorem}[Exact derivative stages from a last exponent]
Write $D_0(x)$ for true, and for $a>0$ set
\[
       D_a(x)\quad\Longleftrightarrow\quad
       x\ne0\ \text{and the last exponent of }x\text{ is at least }a.
\]
For all term indices the following equations hold, where the prime
means taking limit points in the order topology:
\[
 D_{S(a)}=(D_a)',\qquad
 D_\lambda(x)\ \Longleftrightarrow\ (\forall n)\,D_{F(\lambda,n)}(x)
 \ \Longleftrightarrow\
 (\forall a<\lambda)\,D_a(x)\qquad(\lambda\text{ a positive limit}).
\]
Thus a nonzero $x$ with last exponent $e$ belongs at stage $e$ and
disappears at $S(e)$. Zero disappears at stage $1$.
\end{theorem}
\begin{proof}
Consider $x=P\mathbin{\|}[e]$. If $e\le a$, use $(P,S(x))$.
Each point strictly between $P$ and $x$ must extend $P$ with a
first new entry $d<e$. An earlier discrepancy would put it on the
same side of both bounds, and the prefix ending in $e$ already equals
$x$. Its last exponent is at most $d<e\le a$.
When $a=0$ this open interval has no point below $x$. Thus throughout
the case $e\le a$ it contains no other $D_a$ point.
Zero has its own isolating neighborhood.

If $e>a$, fix $b<x$. For $e=S(c)$, the term $F(x,v(b))$ is strictly
between $b$ and $x$ and ends in $c\ge a$. For $e$ a limit, take
$n=\max(v(b),v(a))$. Then $F(x,n)>b$ and its last exponent $F(e,n)$
exceeds $a$, again by Theorem~\ref{thm:fsquant}. This supplies a
$D_a$ point for every prescribed neighborhood and proves the successor
equation.

For the limit equation, $D_\lambda(x)$ implies membership at every
smaller stage by comparison of indices. For nonzero $x$ with last
exponent $e<\lambda$, the one index $F(\lambda,v(e))>e$ witnesses
failure of both universal conditions. For $x=0$, the term
$F(\lambda,0)>0$ witnesses the same failure, since our fundamental
sequences use at least one copy in the bottom clause. These are
explicit witnesses for all the negative cases.
\end{proof}

The family $D$ was defined directly by a primitive recursive test on
two inputs. The displayed equations certify its behavior at every
stage without applying a transfinite recursion principle in iPRA.
Each $D_a$ contains its limit points since $D_{S(a)}\subseteq D_a$.
It is cofinal in the whole term order. For a proposed bound $x$ the
singleton $[S(\max(x,a))]$ exceeds $x$ because $x<[x]$, and its
exponent exceeds $a$. These give explicit closed unbounded families on
this particular countable notation order. The neighborhood proofs also
apply pointwise within $[0,b)$. When $b>0$ the stage $D_b$ in this
segment is empty because, for each positive $x<b$, its last exponent is no
larger than its first, which was proved strictly below $x$.

\section{Controlled descending chains below a fixed finite power of omega}
\label{sec:control}
Fix a positive integer $d$. Terms below $\omega^d=[d]$ are exactly
the polynomials in $\omega$ with natural coefficients and exponents
$d-1,\ldots,0$ (zero coefficients fill the gaps). Indeed, the first
exponent, if present, must be $<d$, and every smaller entry is then
also a finite exponent $<d$. Reading multiplicities from highest to
lowest identifies the term order with strict lexicographic order of
$d$-tuples of naturals. At the first unequal coefficient, the tuple
with the shorter run either ends or passes to a lower exponent, which
proves the comparison rule in both directions.

Let $g(i)$ be a given nondecreasing primitive recursive bound on
all coordinates allowed at time $i$. Define an end index for each
fixed dimension by
\[
 H_0(s)=s+1,\qquad
 H_{j+1}(s)= H_j^{\,g(s)+1}(s),
\]
where the superscript on $H_j$ means ordinary iteration of that
natural-number function. Each $H_j$ is nondecreasing and
$H_j(s)\ge s+1$. Indeed, increasing the
start or the number of iterations cannot decrease the result, and
there is at least one iteration. For every fixed $j$ the iterate
uses only primitive recursion.

\begin{theorem}[A length bound with a computable failure index]
A list occupying times $s,\ldots,t-1$ of strictly decreasing
$j$-tuples satisfying the coordinate bound has $t\le H_j(s)$
(for $s\le t$). In particular every primitive recursive sequence of
terms below $\omega^d$, for fixed $d$, has a primitive recursively
located index $i$ with $f(i+1)\ge f(i)$.
\end{theorem}
\begin{proof}
At dimension zero there is just the empty tuple, so the list has
at most one entry. Suppose the bound is proved in dimension $j$.
In a nonempty decreasing list of $(j+1)$-tuples the first coordinates
are nonincreasing and start at most $g(s)$. Split into maximal
constant-first-coordinate blocks by scanning the list. There are at
most $g(s)+1$ blocks because every new block drops the first coordinate
by at least one. If a block starts at $u$, its remaining $j$ coordinates
decrease strictly, so its exclusive end is at most $H_j(u)$.
Induction on the number of blocks and monotonicity show that the final
end is at most $H_j^{g(s)+1}(s)$. An empty list satisfies the bound
immediately. The dimension induction constructs a separate iPRA proof
for each fixed $j$. Within each proof the block endpoints are bounded
scans of the supplied list, and the induction on block count only
compares natural numbers computed from those endpoints.

For the last assertion choose $g(i)$ as the maximum of the coordinates
of $f(0),\ldots,f(i)$ (and zero). Apply the bound to the list of
$H_d(0)+1$ consecutive inputs starting at zero. If all adjacent
comparisons decrease, its end index violates the bound. A bounded
scan returns a failing comparison before index $H_d(0)$. Its correctness follows by the same finite argument. The statement for
sequences is a scheme, with $f$ given by a primitive recursive
definition. The dimension is fixed before forming the bound.
\end{proof}
